\documentclass[conference]{IEEEtran}
\usepackage[utf8]{inputenc}
\usepackage{amsmath,amssymb}
\usepackage{booktabs}
\usepackage{graphicx}
\usepackage{xcolor}
\usepackage[colorlinks=true,linkcolor=blue,citecolor=blue,urlcolor=blue]{hyperref}

% ============================================================================
% IEEEtran conference format (two-column).
% Every number here is a real measurement. Provenance:
%   corrected closed-loop    -> closed_loop_harness/run_sweep_clean.sh, output/remeasure_clean
%   contaminated closed-loop -> output/benchmark12 (archived, run 2026-06-29)
%   activation crossover     -> icra27/e1 (decoder-only replay; never touches the renderer)
%   manipulation             -> OpenVLA-OFT / LIBERO pipeline (separate harness, NOT yet audited
%                               for the same class of state-reuse bug -- see Limitations)
%
% Double-anonymous: the author block below must stay anonymised for submission.
% ============================================================================

\title{\textbf{Weights Are Free, Activations Are the Wall:\\
Quantizing Driving VLA Policies for Closed-Loop Safety}}

\author{\IEEEauthorblockN{Anonymous Submission}
\IEEEauthorblockA{Paper ID: XXXX}}

\begin{document}
\maketitle


\begin{abstract}
Post-training quantization of vision--language--action (VLA) driving policies is usually validated
on open-loop displacement error, a metric known not to predict closed-loop behaviour. Evaluating
in closed loop instead, we find that \textbf{4-bit weight quantization is free}: on an
impact-speed-graded closed-loop safety benchmark (NeuroNCAP/NeuRAD) over $16$ adversarial scenario
instances $\times\,10$ seeds, naive per-channel INT4 scores $4.71$ ($95\%$ CI $[4.38,5.00]$)
against full precision's $4.81$ $[4.62,4.97]$---a paired difference of $-0.10$, CI
$[-0.33,+0.12]$. INT8 and group-wise INT4 are likewise statistically indistinguishable from FP16.
Plain round-to-nearest at four bits, with no clipping search, codebook, error compensation or
rotation, already matches full precision.

Establishing this required controlling a simulator property we believe is underappreciated. The
render server is \emph{stateful}: an actor-update endpoint mutates the scene, and the evaluation
engine reads the scene once at initialisation. A harness that launches one renderer per scene and
loops quantization configurations inside it therefore gives only the \emph{first} configuration a
clean scene, and because configurations are enumerated in a fixed order the contamination is
perfectly confounded with the treatment. Under that harness we, and an earlier version of this
work, measured an apparent catastrophic 4-bit collapse---graded score $2.62$ against $4.41$ at
FP16---that survived ten seeds, paired comparisons and bootstrap confidence intervals. It was an
artifact; the deficit tracked run \emph{position}, not bit-width. We give the diagnostic
signature, a near-free state-restore fix, and a protocol for closed-loop quantization studies.

The weight axis being free does not make 4-bit deployment free. On a renderer-independent replay
path, quantizing \emph{activations} to four bits with weights left untouched inflates planning
error $47\times$ in-distribution---precisely where weight-only INT4 is lossless---and drives the
rate of degenerate stay-put plans from $42/162$ to $107/162$. A single data-free orthogonal (DCT)
rotation restores full-precision accuracy and returns the degenerate count to base rate. The same
axis and the same fix govern a 7B manipulation policy. The practical conclusion is a division of
labour: spend no effort on weight-grid engineering for driving policies at four bits, and spend it
all on activation outliers. The methodological conclusion is that closed-loop evaluation inherits
the failure modes of the software it runs on, and buys trustworthiness over open-loop metrics only
if simulator state is controlled as carefully as random seeds are.
\end{abstract}
\begin{IEEEkeywords}
vision-language-action models, autonomous driving, post-training quantization, closed-loop
evaluation, reproducibility, simulator state, NeuroNCAP, benchmarking
\end{IEEEkeywords}

\section{Introduction}
Post-training quantization (PTQ) is the standard route to running VLA policies on embedded
compute. How PTQ is \emph{validated} has recently shifted: open-loop, in-distribution accuracy is
widely criticised as a poor proxy for deployed behaviour, and closed-loop simulation is offered as
the remedy. In driving specifically, open-loop displacement error is known not to predict
closed-loop driving quality, and safety-graded closed-loop protocols are increasingly the
reference standard.

This paper is a caution about the remedy. Closed-loop simulators are stateful software systems.
When evaluation harnesses reuse a simulator process across configurations---the natural thing to
do, because simulator startup dominates cost---state from one configuration can leak into the
next. If configurations are enumerated in a fixed order, that leakage is not noise: it is a
systematic bias aligned with the treatment variable, and it is invisible to seed averaging,
paired testing, and confidence intervals, all of which we applied.

We document a concrete instance, quantify it, correct it, and extract a protocol.

\textbf{Contributions.}
\begin{itemize}
\item A \emph{contamination mechanism} in closed-loop VLA evaluation: a stateful render server
combined with fixed-order configuration enumeration, with a diagnostic signature (scene actor
count) that identifies it in existing logs.
\item Evidence that it fabricated a large, stable effect: an apparent $40\%$ loss of graded
closed-loop safety at 4 bits, monotone in \emph{run position} rather than bit-width.
\item A \emph{corrected measurement}: on a restored harness, INT8, naive INT4 and group-wise INT4
are all statistically indistinguishable from FP16 on graded closed-loop safety.
\item A validated, near-free \emph{fix} and an evaluation protocol for closed-loop PTQ studies.
\item A renderer-independent characterisation of what does break at 4 bits---activation
outliers---and evidence that the axis and its fix transfer to a 7B manipulation policy
(Section~\ref{sec:activations}). Together with the corrected weight-axis result, this gives a
division of labour for 4-bit driving deployment: weight grids need no engineering; activation
outliers need a rotation.
\end{itemize}

\section{Setup}
\textbf{Policy.} OpenDriveVLA-0.5B: a Qwen2.5-0.5B planner on UniAD stage-1 track+map perception,
emitting the trajectory as autoregressive text. PTQ is weight-only and applied \emph{only} to the
$168$ backbone Linear layers ($357.8$M parameters); perception, projectors, embeddings and
\texttt{lm\_head} remain in floating point, so any measured effect is attributable to the
language/action head.

\textbf{Closed-loop benchmark.} NeuroNCAP re-renders six photorealistic cameras at the driven ego
pose via NeuRAD, with scripted adversaries in frontal, side and stationary configurations. The
score is impact-speed graded: $5$ if no collision occurs, else
$4\cdot\max(0,\,1-v_\text{impact}/v_\text{ref})$, where $v_\text{ref}$ is the closing speed the ego
would have had on the logged human trajectory. Severity grading matters here: on several scenes
\emph{every} configuration including FP16 collides, so a binary collision rate carries no signal.

\textbf{Aggregation.} $16$ scene/scenario instances $\times\,10$ seeds per configuration. All
aggregates carry $95\%$ percentile-bootstrap confidence intervals over scenario-instance means
($10^4$ resamples); comparisons use paired per-instance differences.

\section{The contamination mechanism}
\label{sec:mech}
The NeuRAD render server exposes \texttt{GET /get\_actors} and \texttt{POST /update\_actors}. The
evaluation engine calls \texttt{get\_actors} \textbf{once at initialisation} and treats the result
as the scene's pristine actor set; each rollout then pushes a scenario-specific actor set via
\texttt{update\_actors}. The push persists in the server.

Consequently the actor set a run observes depends on what ran before it in the same server
process. Measured on scene 0099: a freshly loaded renderer reports $28$ actors; after one
evaluation it reports $1$; a subsequent evaluation initialised from that mutated state logged
$60$ actors, against $33$ in the reference run.

The benchmark harness launched one renderer per \emph{scene} and looped configurations inside it:
\texttt{fp16}, \texttt{naive\_w4}, \texttt{group\_w4\_g128}. FP16 therefore always ran first, on a
clean scene, and each quantized configuration ran on a more heavily mutated one. The extra
background vehicles are placed in the roadway, so contaminated runs collide almost immediately.

\textbf{The effect tracks run position, not bit-width.} Table~\ref{tab:position} shows scene 0099
(stationary). Under contamination the ordering of scores follows the ordering of execution
exactly. On a clean renderer all three configurations are perfect.

\begin{table*}[t]\centering
\caption{Scene 0099 stationary, $10$ seeds. Archived scores follow \emph{execution order};
on a restored renderer the differences vanish.}
\label{tab:position}
\begin{tabular}{clcc}
\toprule
run position & config & archived & restored renderer \\
\midrule
1 (clean) & FP16 & $5.00$ & $5.00$ \\
2 (mutated) & naive W4 & $1.43$ & $\mathbf{5.00}$ \\
3 (mutated twice) & group W4 g128 & $0.00$ & $\mathbf{5.00}$ \\
\bottomrule
\end{tabular}
\end{table*}

This also resolves an anomaly we had previously reported as a scientific finding: group-wise INT4
scored \emph{worse} than naive INT4, despite finer scale granularity being strictly more
expressive. Under the quantization hypothesis this was inexplicable. Under the contamination
hypothesis it is forced---group-wise simply ran third.

\textbf{Diagnostic signature.} The scene actor count logged per rollout identifies the problem in
existing logs without re-running anything: a contaminated run reports a different actor count from
a freshly initialised one ($60$ vs $33$ here). We recommend logging and asserting it.

\section{Corrected measurement}\label{sec:corrected}
\textbf{Fix.} After a fresh renderer launch, capture the pristine actor set once per scene, and
\texttt{POST} it back before \emph{every} evaluation. This is validated: a fresh run and a
restored run both reproduce $5.00$ with $33$ actors on 0099-stationary. Cost is negligible---one
HTTP request---against the $70$--$100$\,s of reloading the renderer. Within a single evaluation
invocation multiple seeds are safe, since all seeds share the initialisation-time snapshot.

\textbf{Result.} Table~\ref{tab:corrected} re-measures the full benchmark under this protocol.
The 4-bit collapse disappears entirely.

\begin{table}[t]\centering
\caption{Graded closed-loop safety, $16$ scenario instances $\times\,10$ seeds, $95\%$ bootstrap CI
over instances. Right column reproduces the contaminated numbers we previously reported.}
\label{tab:corrected}
\begin{tabular}{lccc}
\toprule
config & NCAP (corrected) & $95\%$ CI & previously reported \\
\midrule
FP16 & $4.81$ & $[4.62,\,4.97]$ & $4.41$ \\
W8 (per-channel) & $4.81$ & $[4.58,\,5.00]$ & $4.50$ \\
naive W4 (per-channel) & $\mathbf{4.71}$ & $[4.38,\,5.00]$ & $2.62$ \\
group W4 g128 & $\mathbf{4.86}$ & $[4.67,\,4.98]$ & $2.62$ \\
\bottomrule
\end{tabular}
\end{table}

Paired against FP16 over scenario instances: W8 $\Delta=-0.00$, CI $[-0.15,+0.17]$; naive W4
$\Delta=-0.10$, CI $[-0.33,+0.12]$; group W4 $\Delta=+0.05$, CI $[-0.06,+0.20]$. All three
intervals contain zero. \textbf{Plain round-to-nearest INT4---no clipping search, no codebook, no
error compensation, no rotation---already matches full precision on this benchmark.}

Two honest caveats. First, the corrected scores sit near the metric ceiling ($5.00$ on most
instances), so the benchmark has limited headroom to discriminate \emph{any} recipe; a null result
here is weaker evidence than a null result on a metric with dynamic range. Second, residual
variance is scenario and seed noise rather than quantization: FP16 itself scores $0.00$ on some
seeds of 0108-side and 0346-frontal, exactly as the quantized configurations do.

\section{Weights are free; activations are the wall}
\label{sec:activations}
Section~\ref{sec:corrected} should not be read as ``4-bit driving policies are safe.'' It says the
\emph{weight} axis is not where this policy breaks. The activation axis is, and this section
measures it on a decoder-only replay path that never contacts the renderer and is therefore
structurally immune to the contamination in Section~\ref{sec:mech}: $162$ captured planner inputs
replayed through the planner directly (Table~\ref{tab:act}).

\begin{table*}[t]\centering
\caption{Activation quantization on the driving policy (decoder-only replay of $162$ real inputs;
renderer-independent). FP16 base rate is $42/162$ legitimately-stationary all-zero plans.}
\label{tab:act}
\begin{tabular}{lccc}
\toprule
config & ST-P3 L2 avg (m) & all-zero plans & $|xy|>50$\,m \\
\midrule
FP16 (reference) & $0.595$ & $42$ (base) & $0$ \\
W16A4, no rotation & $\mathbf{27.79}$ ($47\times$) & $107$ & $7$ \\
W16A4 + DCT & $0.646$ & $42$ & $0$ \\
W4A4, no rotation & $9.34$ & $127$ & $3$ \\
\textbf{W4A4 + DCT} & $\mathbf{0.611}$ & $41$ & $0$ \\
\bottomrule
\end{tabular}
\end{table*}

Quantizing activations to $4$ bits with weights untouched inflates planning error $47\times$ and
triples the rate of degenerate stay-put plans, in-distribution, where weight-only INT4 is
lossless. A single data-free orthogonal (DCT-II) rotation, $y=(WQ^\top)(Qx)$, restores FP16-level
accuracy and returns the degenerate-plan count to base rate. The same axis and the same fix govern
a 7B manipulation policy (OpenVLA-OFT on LIBERO): weight-only INT4 is near-lossless
($90.75\%$ vs $90.25\%$ bf16) while naive W4A4 collapses to $0\%$ and rotated W4A4 recovers
$89.8\%$. Activation outliers, not weight granularity, are the binding constraint in both
embodiments.

\section{A protocol for closed-loop PTQ evaluation}
\begin{enumerate}
\item \textbf{Assume the simulator is stateful until proven otherwise.} Any endpoint that mutates
scene contents persists across evaluations sharing a process.
\item \textbf{Snapshot and restore} the mutable state before every evaluation, or relaunch. Assert
the restore succeeded rather than trusting it.
\item \textbf{Randomise or interleave configuration order} across process lifetimes. Fixed
enumeration order converts any residual leakage into a treatment-aligned bias.
\item \textbf{Log an environment invariant} per rollout (here, the scene actor count) and assert
it is constant across configurations. This is the cheapest detector we found and it works
retrospectively on archived logs.
\item \textbf{Treat a full-precision control as a running check}, not a one-off baseline. A
control re-run in a late position would have caught this immediately.
\item \textbf{Distrust monotonicity in enumeration order.} Effects that order themselves the way
the loop does deserve suspicion, especially when they contradict theory---as group-wise scoring
below per-channel did here.
\end{enumerate}

\section{Limitations}
We document one mechanism in one simulator; we do not claim it generalises to CARLA-based or
non-reactive benchmarks, though the preconditions---mutable server state plus process reuse---are
common. The corrected benchmark is near its ceiling, which limits the strength of our null result
and means we cannot bound how much 4-bit damage a more discriminating driving benchmark would
reveal. The activation results are in-distribution replay and are not themselves closed-loop; the
title's second half is therefore a narrower claim than the first---weight-axis freedom is a
closed-loop, safety-graded result, while the activation-axis wall is an accuracy result on a
replay path. The manipulation cross-embodiment evidence (Section~\ref{sec:activations}) comes from
a separate LIBERO harness that we have \emph{not} audited for the class of state-reuse bug
identified in Section~\ref{sec:mech}; given that a nominally identical bug produced a stable, false
effect in our own driving benchmark, we flag this as an open risk rather than assume the
manipulation harness is exempt, and we do not present it as independently re-verified. Each
embodiment is a single checkpoint. Finally, closed-loop success throughout is measured with
fake-quant weights, so we report no measured latency win.

\section{Conclusion}
We set out to find a 4-bit quantization recipe that would recover closed-loop driving safety lost
at 4 bits. There was no loss to recover. The deficit we and our earlier draft had measured was
produced by a stateful simulator combined with fixed-order configuration enumeration, and it
withstood every statistical control we applied because those controls address noise, not bias
aligned with the treatment. On a corrected harness, plain INT4 matches full precision. We
therefore offer the negative result as the contribution: closed-loop evaluation inherits the
failure modes of the software it runs on, and the field's current move from open-loop to
closed-loop metrics buys trustworthiness only if simulator state is controlled as carefully as
seeds are.

\section*{Reproducibility}
The restore-protocol harness, the aggregation and diagnostic scripts, the contaminated and
corrected result trees, and the analysis that identified the mechanism are released with the
paper.

\section*{Acknowledgments}
% ICRA requires any AI-generated content to be disclosed in the acknowledgments; ICLR instead
% requires a dedicated (non-page-counted) AI-usage section. Complete per the target venue.
\emph{Disclosure of AI assistance to be completed per target-venue policy.}

\end{document}
